% =====================================================================================================================
\section{Discussion}
\label{sec:discussion}

Software engineering settled long ago where checks belong: a type error is caught when the code is compiled, not by
the colleague who reads the pull request. The numbers in an agent-written paper are in the same position. The reviewer
in the loop reads the paper and even names the wrong value, but its score does not carry that judgment back to the
writer. An audit of the finished paper can find errors, and a strong one finds most of them, but it reports after the
agent has moved on. A recipe rerun at the edit closes that loop where the error is made, in the same turn, with the
correct value attached, and our placement study shows that an agent who receives that message fixes the number instead
of defending it.

The result is less about our checker than about where any checker sits. Agent harnesses expose hooks around tool
calls, and most of what is written about them concerns safety: blocking a destructive command or a write outside the
workspace. The same hooks are a place to run domain checks that are cheap, deterministic, and specific, and our
measurements say that this is where such checks pay off. Telling an agent in its prompt that a check exists is not
the same as running it. The prompted tool and the automatic hook contain the same code and the same rules. For our main
writer, one leaves the paper as it would be with no check and the other removes every planted error. Other writers
left fewer planted errors with the offer, and the largest left none, but a user cannot count on which writer they
have; the hook gave the same outcome for all three.

The recipe also changes what a paper carries. Each number arrives with the computation that produced it and a
sentence saying what that computation means, which a co-author can read, a reviewer can rerun, and a second model can
compare against the prose. Generate mode goes one step further and removes the need to type a result number by hand.
The two checkers we studied fail on different numbers, which suggests a division of labor: recomputation at every
edit for values and rounding, and an occasional strong audit for meaning.

\paragraph{Limitations and open directions.}
The placement study covers 8 repositories and three writer models from one family, run in one agent harness, and we
ran the arm without any check for only one of the three. Its planted
errors sit on numbers that already carry a working recipe, which is the favorable case for a recompute check: an
error whose recipe is itself wrong belongs to the blind spot measured in Section~\ref{sec:checkers}. The pre-commit
path brings the check to other agents and editors, and open-weight writers are a natural next test. The check trusts
the writer's statement of what a number means, and a check of the \code{means} string against the sentence would
close the gap behind both of recompute's misses. Our instruction to protect the results is one
sentence in a prompt. A writer that is trained or rewarded to keep a result is a harder test, and so are sessions in
which the agent makes no edit and only the end-of-session check fires; blocking could still matter in both. Settings and conventions such as confidence
levels need an allowed list, and numbers the agent never tags need a coverage report. Our reviewer study measures
scores on a percentile scale; reviewers that are trained or prompted to act on the errors they find may behave
differently, and a reviewer that receives the hook's log beside the paper is a design we have not tried.
